\chapter{Curvas Simples. Clasificación. Teorema de la Curva de Jordan. Teorema de la Divergencia}

\begin{definicion}
    Una \textbf{curva simple} es un subconjunto $C\subset \bb{R}^2$ no vacío tal que $\forall p\in C$ se tiene que 
    \begin{itemize}
        \item $\exists\ V$ entorno abierto de $p$ en $C$.
        \item $\exists\ I\subset \bb{R}$ intervalo abierto.
        \item $\exists\ \alpha:I \to \bb{R}^2$ curva parametrizada regular
    \end{itemize} 
    de forma que $\alpha(I)=V$, y la aplicación $\alpha:I\to V$ es un homeomorfismo.
\end{definicion}

% dibujo

% En cierta forma estamos copiando la definición de superficies pero para R^2.

\begin{observacion}\
    \begin{enumerate}
        \item La aplicación $\alpha$ la llamaremos \textbf{parametrización local} de $C$.
        \item $C$ no puede tener ``esquinas'' (o ``picos'') ya que $\alpha$ es regular.
        \item $C$ no puede tener autointersecciones. Ni siquiera el Folium de Descartes puede ser una curva simple, a pesar de que en este caso $\alpha$ es inyectiva.
        \item Los entornos $V$ son todos de la forma
        \begin{gather*}
            V = U\cap C
        \end{gather*}
        donde $U$ es un entorno abierto de $p$ en $\bb{R}^2$.
        \item Los entornos $V$ son todos conexos.
    \end{enumerate}
\end{observacion}

\begin{ejemplo}
    Supongamos que tenemos una curva simple $C\subset \bb{R}^2$ y consideramos las parametrizaciones $\alpha: I \to C$ y $\beta:J \to C$ de forma que $W = \alpha(I)\cap \beta(J)\neq \emptyset$.\\

    Tenemos que $\alpha^{-1}(W)\subset I$ y que $\beta^{-1}(W)=J$. Además, como $W$ es abiertos, tendremos que $\alpha^{-1}(W)$ es abierto en $I$ y $\beta^{-1}(W)$ es abierto en $J$ y por ser abiertos de un abierto, ambos serán abiertos en $\bb{R}$.\\

    Además, podemos considerar
    \begin{gather*}
        \varphi = \beta^{-1}\circ \alpha : \alpha^{-1}(W) \to \beta^{-1}(W)\\
        \psi = \alpha^{-1}\circ \beta : \beta^{-1}(W) \to \circ^{-1}(W)
    \end{gather*}
    Veamos que $\varphi$ y $\psi$ son \textbf{difeomorfismos}.
\end{ejemplo}